\documentclass[11pt]{article}
\usepackage[a4paper,margin=1in]{geometry}
\usepackage{amsmath,amssymb,mathtools,bm}
\usepackage{enumitem,booktabs,array}
\usepackage{tcolorbox}
\usepackage{hyperref}
\usepackage[protrusion=true,expansion=false]{microtype}
\hypersetup{colorlinks=true,linkcolor=blue,urlcolor=blue,citecolor=blue}
\setlist[itemize]{topsep=2pt,itemsep=2pt}
\setlist[enumerate]{topsep=2pt,itemsep=2pt}
\newtcolorbox{keybox}{colback=blue!3!white,colframe=blue!40!black,boxrule=0.4pt,arc=1mm}
\newtcolorbox{alertbox}{colback=red!3!white,colframe=red!40!black,boxrule=0.4pt,arc=1mm}
\newtcolorbox{answerbox}{colback=green!3!white,colframe=green!40!black,boxrule=0.4pt,arc=1mm}
\newtcolorbox{drillbox}{colback=yellow!5!white,colframe=orange!60!black,boxrule=0.4pt,arc=1mm}
\newcommand{\EV}{\mathbb{E}}
\newcommand{\Var}{\mathrm{Var}}
\newcommand{\Cov}{\mathrm{Cov}}
\newcommand{\blank}[1]{\underline{\hspace{#1}}}

\title{W05-01: Iyer \& Puri (2012, AER) \\
\large ``Understanding Bank Runs: The Importance of Depositor-Bank Relationships and Networks'' --- Active Recall Workbook}
\author{Banking \& Risk Management --- Exam Prep}
\date{\today}
\begin{document}\maketitle

\begin{keybox}
\textbf{Paper in one sentence.} Using minute-by-minute depositor-level data from a 2001 run on a failing Indian cooperative bank, IP show that (i) insured depositors still run, (ii) longer relationships and depositor loans \emph{reduce} run probability, and (iii) social network links (same ethnicity, same branch, introducer ties) transmit the run --- modern empirical backing of Diamond--Dybvig, with behavioural and network dimensions absent in the theory.

\textbf{Placement.} The cleanest depositor-level run study. Complements Gorton--Metrick (W04-01) on wholesale runs with a \emph{retail} run story and adds a network/behavioural angle.
\end{keybox}

\section*{Round 1: Complete Walk-Through}

\subsection*{1.1 The bank and the shock}
A cooperative bank in India (``Bank A'') fails in early 2001 amid fraud revelations at a peer bank. On run days, depositors queue to withdraw. Over subsequent weeks partial withdrawal restrictions apply. Deposits are partly insured up to 100{,}000 INR (then about \$2{,}200) by the DICGC.

\subsection*{1.2 Data}
Minute-level withdrawal records: account-holder ID, amount withdrawn, branch, demographic info, relationship length, loan status, same-branch introducer. IP hand-collect a depositor-to-depositor introducer network and proxy ethnicity.

\subsection*{1.3 Central findings}
\begin{itemize}
\item \textbf{Deposit insurance is not a full firewall.} Fully insured depositors \emph{also} run, although less than uninsured.
\item \textbf{Relationships matter.} Longer account tenure reduces run probability $\Rightarrow$ ``soft information'' and switching costs help.
\item \textbf{Loan relationships.} Depositors who \emph{also borrow} from the bank run much less --- another friction against running.
\item \textbf{Network contagion.} A depositor is more likely to run if their \emph{introducer} ran, or if depositors at their branch in the same ethnic group ran: social transmission of the run is econometrically identifiable.
\end{itemize}

\subsection*{1.4 Estimation}
Hazard / linear probability models:
\[
\text{Run}_{it}=\alpha + \beta_1\cdot\text{Insured}_i+\beta_2\cdot\text{Tenure}_i+\beta_3\cdot\text{Loan}_i+\beta_4\cdot\text{IntroducerRan}_{it}+X_i'\gamma+\varepsilon_{it}.
\]
$\beta_4>0$: depositors imitate their referrer. Time-dependent covariates (``introducer ran in the last 30 minutes'') rule out common-shock stories.

\subsection*{1.5 Caveats}
\begin{alertbox}
\begin{itemize}
\item Single-bank study; external validity is limited to similar retail settings.
\item Network measures are noisy (introducer information is self-reported at account opening).
\item Social transmission could partly reflect common information arriving via friends, not pure contagion; IP discuss this carefully.
\end{itemize}
\end{alertbox}

\section*{Round 2: Level 1 Fill-in}
\begin{enumerate}
\item Setting: \blank{1cm} cooperative bank, year \blank{1cm}.
\item Data granularity: \blank{1cm}-level withdrawals.
\item Relationship protection: longer \blank{1.5cm} and having a \blank{1cm} reduce running.
\item Network variable: depositor's \blank{2cm} ran.
\item Insured depositors: do \blank{1cm} run (yes/no).
\end{enumerate}

\section*{Round 3: Level 2 Prompts}
\begin{enumerate}
\item Write a simple run-hazard specification and state what each covariate identifies.
\item Why do insured depositors run? What does this say about the micro foundations of deposit insurance?
\item How do IP disentangle common-information shocks from pure social contagion?
\item Link IP to Diamond--Dybvig; what dimensions does IP add?
\end{enumerate}

\section*{Round 4: Minimal Scaffolding}
From a blank page: (i) setting and data, (ii) four central findings, (iii) hazard spec, (iv) network-contagion story, (v) two caveats.

\section*{Round 5: Blank Page}
Essay: what do the Indian retail-run data teach about modern bank-run models?

\vspace{6pt}
\begin{drillbox}
\textbf{Speed Drill.} (1) Country/year. (2) Granularity. (3) Two ``relationship'' protections. (4) Network variable. (5) Fact about insured depositors.
\end{drillbox}

\section*{Quick Reference \& Answer Key}
\begin{answerbox}
\begin{itemize}
\item \textbf{Setting.} Indian cooperative bank, 2001 failure.
\item \textbf{Relationships.} Longer tenure and loan ties cut run probability.
\item \textbf{Insurance.} Not a perfect firewall --- insured depositors also run.
\item \textbf{Networks.} Introducer-ran and same-ethnicity-ran effects: social transmission.
\end{itemize}
\end{answerbox}

\begin{keybox}
\textbf{30-second exam pitch.} Iyer \& Puri (2012) provide the cleanest depositor-level evidence on bank runs. Minute-by-minute data from an Indian cooperative bank that failed in 2001 shows that (i) even fully insured depositors run --- deposit insurance is an imperfect firewall; (ii) longer account tenure and having a loan with the bank reduce run probability, reflecting relationship-based soft-information frictions; (iii) social networks transmit the run --- depositors are more likely to withdraw after their introducer or same-ethnicity neighbours withdraw. The paper puts empirical teeth on Diamond--Dybvig at the retail level and integrates behavioural and social-network amplification into the bank-run literature. Policy message: insurance alone does not stop runs; relationships and community ties shape both the run's spread and its containment.
\end{keybox}

\end{document}
